\section*{Chapitre \ref{ch:b2:current-potential-curves} --- Courbes intensité--potentiel}

\begin{solution}{exo:b2:current-potential-curves:1}
Le zinc est oxydé~: \omterm{def:b2:current-potential-curves:ie-curve}{courant anodique}, positif. Les ions cuivre sont réduits à
l'électrode de cuivre~: \omterm{def:b2:current-potential-curves:ie-curve}{courant cathodique}, négatif. Les deux courants sont égaux en
valeur absolue, puisque la même charge traverse le circuit.
\end{solution}

\begin{solution}{exo:b2:current-potential-curves:2}
$|i_{\lim}| = 96\,485 \times 0.50 \times 10^{-4} \times 1.6 \times 10^{-9} \times
2.0/(25 \times 10^{-6}) = \qty{6.2e-4}{A}$, soit \qty{0.62}{mA}.
\end{solution}

\begin{solution}{exo:b2:current-potential-curves:3}
Un \omterm{def:b2:current-potential-curves:fast-slow}{système rapide} coupe l'axe des potentiels brutalement en son potentiel de Nernst~; un \omterm{def:b2:current-potential-curves:fast-slow}{système
lent} présente autour de lui une portion plate à courant nul, ses branches ne commençant
qu'après une \omterm{def:b2:current-potential-curves:overpotential}{surtension}. Sur platine~: \ce{Fe^3+}/\ce{Fe^2+} (rapide)~;
\ce{O2}/\ce{H2O} (lent).
\end{solution}

\begin{solution}{exo:b2:current-potential-curves:4}
Au potentiel de demi-vague, égal à $E^\circ = \qty{0.77}{V}$ lorsque les deux formes
diffusent de la même façon.
\end{solution}

\begin{solution}{exo:b2:current-potential-curves:5}
Vers \qty{0.77}{V}, une vague cathodique de \ce{Fe^3+} (rapide), qui atteint un palier
proportionnel à sa concentration. Vers \qty{0.34}{V}, une seconde vague, due au dépôt du
cuivre~; son palier s'ajoute au premier et est environ deux fois plus haut pour une concentration
égale et des coefficients de diffusion voisins, puisque $n = 2$. Puis, vers
\qty{0}{V} à pH 0 sur platine (plus bas une fois l'électrode recouverte de cuivre),
le mur abrupt de la réduction de \ce{H+}.
\end{solution}

\begin{solution}{exo:b2:current-potential-curves:6}
À \qty{0.50}{V}~: seul \ce{Fe^3+} est réduit en \ce{Fe^2+}, à son \omterm{def:b2:current-potential-curves:limiting-current}{courant limite}.
À \qty{0.10}{V}~: \ce{Fe^3+} est réduit et le cuivre se dépose simultanément,
chacun à son \omterm{def:b2:current-potential-curves:limiting-current}{courant limite}~; le total est la somme des deux paliers.
\end{solution}

\begin{solution}{exo:b2:current-potential-curves:7}
pH 0~: \qty{0}{V} et \qty{1.23}{V}~; pH 14~: \qty{-0.83}{V} et \qty{0.40}{V}. Sur
platine, le domaine garde sa largeur et descend de \qty{59}{mV} par unité de pH,
chaque mur suivant son couple, le mur anodique restant décalé vers le haut par la \omterm{def:b2:current-potential-curves:overpotential}{surtension}
du dioxygène.
\end{solution}

\begin{solution}{exo:b2:current-potential-curves:8}
À \qty{1.0}{V}, \ce{Fe^2+} est oxydé sur son palier (\omterm{def:b2:current-potential-curves:ie-curve}{courant anodique}
proportionnel à $[\ce{Fe^2+}]$) et tout \ce{Ce^4+} est réduit sur son palier
(\omterm{def:b2:current-potential-curves:ie-curve}{courant cathodique} proportionnel à $[\ce{Ce^4+}]$). Avant l'équivalence, le
\omterm{def:b2:current-potential-curves:ie-curve}{courant anodique} décroît linéairement avec le volume ajouté~; après, le \omterm{def:b2:current-potential-curves:ie-curve}{courant cathodique}
croît linéairement. Les deux droites coupent le zéro à l'équivalence.
\end{solution}

\begin{solution}{exo:b2:current-potential-curves:9}
Les paliers doublent ($|i_{\lim}| \propto 1/\delta$). Le potentiel de demi-vague ne
change pas, la couche étant la même pour les deux formes~; le potentiel à courant nul
d'un \omterm{def:b2:current-potential-curves:fast-slow}{système rapide}, qui est le potentiel de Nernst, ne change pas non plus.
\end{solution}

\begin{solution}{exo:b2:current-potential-curves:10}
Comme dans le chapitre avec $i_{\lim,c} = 0$~: $E = E_{1/2} + (RT/nF)\ln[i/(i_{\lim} -
i)]$. En logarithmes décimaux, $E = E_{1/2} + (RT\ln 10/nF)\log[i/(i_{\lim} - i)]$,
une droite de pente $0.0592/n$ volt à \qty{298}{K}, qui coupe l'axe des
$E$ en $E_{1/2}$.
\end{solution}

\begin{solution}{exo:b2:current-potential-curves:11}
Seul, le zinc se fixe au potentiel où son \omterm{def:b2:current-potential-curves:ie-curve}{courant anodique} (oxydation rapide
du zinc) est égal au \omterm{def:b2:current-potential-curves:ie-curve}{courant cathodique} de réduction de \ce{H+} sur le zinc,
qui est faible parce que cette réduction est lente sur zinc~: la dissolution est
lente. Au contact du platine, les électrons libérés par le zinc peuvent réduire \ce{H+}
sur le platine, où la réduction est rapide~: le \omterm{def:b2:current-potential-curves:ie-curve}{courant cathodique} à un potentiel donné
est bien plus grand, l'équilibre des courants est atteint à un courant plus élevé, et le
zinc se dissout plus vite pendant que du dihydrogène se forme sur le platine.
\end{solution}

\begin{solution}{exo:b2:current-potential-curves:12}
Sur platine, l'oxydation de l'eau commence vers \qty{1.6}{V} à pH 0~: avant
\qty{2.0}{V}, le courant servirait à produire du dioxygène. Une électrode sur laquelle
l'oxydation de l'eau exige une \omterm{def:b2:current-potential-curves:overpotential}{surtension} plus grande repousse le mur au-delà de
\qty{2.0}{V}, de sorte que l'espèce peut être oxydée à l'intérieur du \omterm{def:b2:current-potential-curves:window}{domaine d'électroactivité}.
\end{solution}

\begin{solution}{pb:b2:current-potential-curves:1}
\textbf{1.} \ce{O2 + 2H2O + 4e- -> 4OH-}.
\textbf{2.} \ce{Ag + Cl- -> AgCl + e-} (quatre fois)~; réaction globale \ce{O2 + 2H2O + 4Ag +
4Cl- -> 4AgCl + 4OH-}.
\textbf{3.} $E = 1.229 - 0.0592 \times 7 + (0.0592/4)\log 0.21 = \qty{0.80}{V}$.
\textbf{4.} $E^\circ(\ce{AgCl}/\ce{Ag}) = \qty{0.22}{V}$ avec des ions chlorure d'activité 1~:
$-0.60 + 0.22 = \qty{-0.38}{V}$ sur l'échelle de l'hydrogène.
\textbf{5.} La réduction de \ce{O2} est un \omterm{def:b2:current-potential-curves:fast-slow}{système lent}~: au-dessous de son potentiel de Nernst,
il faut une \omterm{def:b2:current-potential-curves:overpotential}{surtension} de plusieurs dixièmes de volt avant que le courant
augmente, et davantage pour atteindre le palier.
\textbf{6.} Sur le palier, le courant ne dépend que de l'apport de \ce{O2},
et non de petites dérives du potentiel~: il mesure la concentration.
\textbf{7.} La réduction de l'eau en dihydrogène (le mur cathodique), qui
ajouterait un courant sans rapport avec le dioxygène.
\textbf{8.} Linéaire à travers la membrane, de $c$ côté échantillon à zéro à
la cathode.
\textbf{9.} $j = D_m c/L$ par unité de surface.
\textbf{10.} $i = -4FAD_m c/L$ (\omterm{def:b2:current-potential-curves:ie-curve}{courant cathodique}), $|i| = 4FAD_mc/L$.
\textbf{11.} La membrane est l'étape la plus lente~: la diffusion à travers elle est bien
plus lente que dans l'eau extérieure, si bien que la membrane joue le rôle de la
\omterm{def:b2:current-potential-curves:limiting-current}{couche de diffusion}.
\textbf{12.} Le gradient de concentration se situe dans la membrane, dont l'épaisseur
ne dépend pas de l'agitation (pourvu que l'eau extérieure soit renouvelée)~;
un dépôt sur la membrane en augmente l'épaisseur et abaisse le courant.
\textbf{13.} $0.2095 \times (101.325 - 3.17) = \qty{20.56}{kPa}$.
\textbf{14.} $c = 1.3 \times 10^{-5} \times 20\,560 = \qty{0.27}{mol/m^3}$ ($0.267$).
\textbf{15.} $0.267 \times 32.00 = \qty{8.6}{mg/L}$ (\unit{g/m^3}).
\textbf{16.} $|i| = 4 \times 96\,485 \times 1.0 \times 10^{-5} \times 1.0 \times 10^{-10}
\times 0.267/(25 \times 10^{-6}) = \qty{4.1}{\micro A}$.
\textbf{17.} $4.05/8.55 = \qty{0.474}{\micro A}$ par \unit{mg/L}.
\textbf{18.} L'épaisseur de la membrane et son coefficient de diffusion ne sont pas connus
avec précision et changent avec l'âge et la température~; une mesure dans une solution
connue les intègre tous.
\textbf{19.} Zéro (en pratique un faible courant résiduel, que l'on retranche).
\textbf{20.} $2.80/0.474 = \qty{5.9}{mg/L}$.
\textbf{21.} $2.80/4.05 = 69~\%$ de la saturation.
\textbf{22.} Du dioxygène consommé par la décomposition de matière organique (eaux usées, végétaux morts)
plus vite qu'il n'est apporté par l'air et par la photosynthèse.
\textbf{23.} $2.80 \times 10^{-6}/(4 \times 96\,485) \times 3600 = \qty{2.6e-8}{mol}$
par heure~: négligeable devant le dioxygène contenu dans un seul litre d'eau
(\qty{1.8e-4}{mol}).
\textbf{24.} La solubilité de \ce{O2} (\omterm{prop:b2:chemical-potential:henry}{constante de Henry}) et la diffusion
à travers la membrane changent toutes deux avec la température~: la sensibilité déterminée à
\qty{25}{\degreeCelsius} ne s'applique plus.
\textbf{25.} L'échantillon contient $\boldsymbol{\approx \qty{5.9}{mg/L}}$ de dioxygène
dissous.
\end{solution}
